\documentclass[12pt]{jlreq}
\pagestyle{empty}
\usepackage[default]{fontsetup}
\usepackage{derivative,fixdif}
\DeclareMathOperator{\atantwo}{atan2}
\newcommand{\bm}[1]{\symbfit{#1}}
\setlength{\parindent}{0pt}
\begin{document}
\section{2次元直交座標}
\begin{gather*}
    \left\{
        \begin{aligned}
            \bm{e}_x&=\left(1,0\right)\\
            \bm{e}_y&=\left(0,1\right)
        \end{aligned}
    \right.
\end{gather*}
\(\bm{e}_x,\bm{e}_y\)は\(t\)に対して不変(定数のように扱われる)\\
\begin{align*}
\bm{r}(t)&=x(t)\bm{e}_x+y(t)\bm{e}_{y}\\
\odv{\bm{r}(t)}{t}&=\odv{x(t)}{t}\bm{e}_x+\odv{y(t)}{t}\bm{e}_{y}\\
\odv[2]{\bm{r}(t)}{t}&=\odv[2]{x(t)}{t}\bm{e}_x+\odv[2]{y(t)}{t}\bm{e}_{y}
\end{align*}
\section{2次元極座標}
基本の変換
\begin{gather*}
    \left\{
        \begin{aligned}
            x(t)&=r(t)\cos\theta(t)\\
            y(t)&=r(t)\sin\theta(t)
        \end{aligned}
    \right.\\
    \left\{
        \begin{aligned}
            r(t)&=\sqrt{x^2(t)+y^2(t)}\\
            \theta(t)&=\atantwo\left(y(t),x(t)\right)
        \end{aligned}
    \right.
\end{gather*}
\subsection{基底ベクトルの変換}
直交座標からの変数変換
\[
(x_1(t),\dots,x_n(t))\to(x_1((y_1(t),\dots,y_n(t))),\dots,x_n((y_1(t),\dots,y_n(t))))
\]
により\[\bm{r}(y_1(t),\dots,y_n(t))=\displaystyle\sum_{k=1}^{n}x_k(y_1(t),\dots,y_n(t))\bm{e}_{x_k}\]として表現可能な場合、各変数方向の基底ベクトルは
\begin{equation*}
    \bm{e}_{y_k}(y_1(t),\dots,y_n(t))=\frac{\pdif{y_k}\bm{r}(y_1(t),\dots,y_n(t))}{\left|\pdif{y_k}\bm{r}(y_1(t),\dots,y_n(t))\right|}
\end{equation*}
で表される。
\end{document}